\chapter[Material extendido]{Material extendido de diseño, desarrollo y evaluación}
\label{ap:material-extendido}

Este anexo conserva el detalle que se retiró del cuerpo principal para cumplir
el límite de páginas. El cuerpo mantiene la línea argumental y las decisiones
centrales. Este anexo permite revisar el razonamiento completo sin cargar los
capítulos de diseño, desarrollo y pruebas.

\section{Diseño extendido del pipeline}
\label{ap:diseno-extendido}

El problema de diseño no consiste solo en generar código. El sistema debe
transformar conocimiento científico en comportamiento simulado. Esa
transformación exige varios cambios de representación: investigación,
formulación matemática, adaptación al entorno, generación de código y
validación.

Cada cambio reduce ambigüedad. La investigación decide qué marcos científicos
son relevantes. La formalización convierte esos marcos en variables, parámetros
y ecuaciones. El Reasoner adapta esas ecuaciones a acciones y percepciones
disponibles. El Builder traduce la especificación a código ejecutable y tests.
Separar estas etapas permite saber si un fallo procede de una fuente
bibliográfica, de una formulación incoherente, de una adaptación imposible o de
un error de implementación.

\subsection{Router y control de estados}

El Router coordina el pipeline mediante una máquina de estados explícita. No se
dejó esta decisión a un planificador conversacional, porque la ejecución tenía
que ser reanudable, auditable y compatible con revisión humana.

El Router conserva el estado del pipeline. Ese estado incluye la etapa actual,
los paradigmas aprobados, las formulaciones seleccionadas, las specs aprobadas
y los artefactos registrados. También emite eventos de traza para reconstruir
herramientas llamadas, tiempos, errores y artefactos producidos.

Las etapas de memoria son condicionales. Se ejecutan cuando Neo4j, Qdrant y los
servicios de embeddings están disponibles. Si esos servicios fallan, el
pipeline puede seguir generando modelos en modo degradado. En cambio,
PostgreSQL y MinIO son centrales: sin ellos no hay persistencia fiable de
estado ni de artefactos.

\begin{table}[H]
\begin{center}
\small
\begin{tabular}{|p{3.1cm}|p{4.8cm}|p{4.6cm}|} \hline
\textbf{Decisión} & \textbf{Motivo} & \textbf{Efecto en el sistema} \\ \hline
Máquina de estados & Evitar que el orden de ejecución dependa de una respuesta LLM. & Reanudación, depuración y trazabilidad. \\ \hline
Estado persistido & No confiar en la memoria conversacional. & Recuperación tras interrupciones. \\ \hline
Rutas de corrección & Volver solo a la etapa responsable del error. & Menos coste y menos repetición innecesaria. \\ \hline
Memoria condicional & No bloquear la generación por servicios opcionales. & Degradación controlada. \\ \hline
\end{tabular}
\caption{Motivos del diseño del Router.}
\end{center}
\end{table}

\subsection{Agentes y subagentes}

El sistema usa agentes principales y subagentes porque cada etapa tiene una
unidad natural de trabajo. Un paradigma puede investigarse sin cargar todos los
demás. Una formulación puede adaptarse sin mezclar su contexto con otras. Una
spec puede implementarse sin leer todo el historial del pipeline.

\begin{table}[H]
\begin{center}
\small
\begin{tabular}{|p{3.2cm}|p{4.6cm}|p{4.7cm}|} \hline
\textbf{Etapa} & \textbf{Agente principal} & \textbf{Subagentes} \\ \hline
Researcher & Descubre candidatos y sintetiza el informe global. & DeepResearchers, uno por paradigma. \\ \hline
Formalizer & Reparte paradigmas aprobados y registra formulaciones. & Un subagente por paradigma. \\ \hline
Reasoner & Empareja formulaciones con el entorno. & Un subagente por paradigma y selección. \\ \hline
Builder & Reparte specs aprobadas y registra modelos. & Un subagente por spec implementable. \\ \hline
\end{tabular}
\caption{Patrón de agente principal y subagentes.}
\end{center}
\end{table}

Esta división controla la ventana de contexto. El agente padre conserva
resúmenes y rutas de artefactos. Cada subagente recibe solo los ficheros que
necesita. Esto reduce interferencias entre paradigmas y evita que variables,
autores o supuestos de una rama aparezcan en otra.

También mejora el aislamiento de fallos. Si una formulación no es válida, no
invalida todo el run. Si un modelo falla tests, el sistema puede repetir solo el
Builder para esa especificación o volver al Reasoner si el problema está en la
adaptación.

\subsection{Researcher y DeepResearcher}

El Researcher mantiene la visión científica de conjunto. Primero consulta la
memoria para recuperar paradigmas conocidos, nombres alternativos y relaciones
previas. Después usa búsqueda acotada para descubrir huecos reales en la
cobertura.

Cada paradigma candidato se convierte en una tarea de investigación profunda.
El DeepResearcher tiene un objetivo único: investigar ese paradigma con su
propia ventana de contexto. Produce un informe con fundamentos, postulados,
supuestos, predicciones, variables y referencias.

El Researcher no carga todos los informes completos en su contexto. Recibe
resúmenes, conserva rutas y sintetiza el informe global. Esta jerarquía separa
cobertura y profundidad.

Los artefactos de esta etapa tienen forma estable:

\begin{verbatim}
research/{run_id}/report.md
research/{run_id}/deep/{paradigm}.md
\end{verbatim}

Antes de emitir la salida estructurada, el Researcher consulta slugs conocidos.
La emisión final queda restringida a esos slugs o a una marca de nuevo
paradigma. Si aparece un paradigma nuevo, el sistema lo canonicaliza antes de
materializarlo en memoria. Esto evita fragmentar la base de conocimiento con
nombres casi equivalentes.

\subsection{Formalizer}

El Formalizer recibe los paradigmas aprobados y genera formulaciones
matemáticas. El componente principal actúa como orquestador: lanza un subagente
por paradigma aprobado y registra los artefactos.

Cada subagente lee el informe profundo de su paradigma y escribe un documento
de formulaciones. El contexto se limita a ese paradigma para evitar que una
ecuación, variable o supuesto se mezcle con otra línea científica.

Cada formulación debe incluir variables, parámetros, ecuaciones y lógica de
decisión. Las formulaciones de un mismo paradigma deben ser estructuralmente
distintas. No basta con renombrar parámetros.

El artefacto principal es:

\begin{verbatim}
research/{run_id}/formulations/{paradigm}.md
\end{verbatim}

Después de esta etapa, el usuario selecciona qué formulaciones continúan. Esta
decisión evita que variantes débiles o redundantes pasen a Reasoner.

\subsection{Reasoner}

El Reasoner adapta las formulaciones seleccionadas al entorno de simulación. Su
entrada incluye la formulación y el archivo \texttt{env\_spec.json}, que define
acciones, percepción, recursos, recompensas y restricciones.

Cada subagente lee:

\begin{verbatim}
deep/{paradigm}.md
formulations/{paradigm}.md
env_spec.json
\end{verbatim}

y produce una especificación JSON por formulación válida:

\begin{verbatim}
models/{run_id}/reasoner/{paradigm}/{formulation}.json
\end{verbatim}

La spec JSON es el contrato entre Reasoner y Builder. Contiene variables,
parámetros, reglas de actualización, lógica de decisión, mapeo al entorno,
comportamientos esperados y referencias.

Antes de escribir una spec, el Reasoner valida que las acciones existan, que las
variables estén definidas, que no haya dependencias circulares incoherentes y
que las percepciones usadas sean observables o computables.

Si una formulación no se puede adaptar, el Reasoner no fuerza una
implementación. Escribe una salida con \texttt{status: invalid} y problemas
explícitos. Así el fallo queda documentado en el nivel correcto.

\subsection{Builder}

El Builder traduce specs aprobadas a modelos Python y tests. El Builder
principal reparte una tarea por especificación. Cada subagente genera:

\begin{verbatim}
models/{run_id}/builder/{paradigm}/{formulation}_model.py
models/{run_id}/builder/{paradigm}/test_{formulation}.py
\end{verbatim}

El modelo debe ser autocontenido y cumplir el contrato:

\begin{lstlisting}[language=Python]
def decide(self, perception: dict) -> Action
def update(self, action, reward, new_perception) -> None
def get_state(self) -> dict
\end{lstlisting}

El Builder ejecuta \texttt{uv run pytest}. Si las pruebas fallan, el subagente
puede corregir el código y repetir hasta un límite de intentos. Si la spec usa
claves de percepción inexistentes o comportamientos no testeables, guarda una
validación en vez de producir código defectuoso.

\subsection{Herramientas de los agentes}

Los agentes no comparten todas las herramientas. Cada etapa recibe solo las que
necesita para su responsabilidad.

\begin{table}[H]
\begin{center}
\small
\begin{tabular}{|p{2.8cm}|p{4.2cm}|p{5.2cm}|} \hline
\textbf{Agente} & \textbf{Herramientas} & \textbf{Uso} \\ \hline
Researcher & \texttt{retrieve\_knowledge}, búsqueda web, investigación profunda, lectura de informes. & Reutilizar memoria, detectar huecos y sintetizar cobertura. \\ \hline
DeepResearcher & Búsqueda académica, búsqueda web y memoria cuando procede. & Investigar un paradigma con contexto propio. \\ \hline
Formalizer & Lectura, escritura y memoria. & Convertir investigación aceptada en formulaciones. \\ \hline
Reasoner & Lectura, escritura y memoria. & Adaptar formulaciones al entorno sin inventar observables. \\ \hline
Builder & Lectura, escritura, \texttt{run\_tests} y memoria. & Generar código, tests y validar ejecución. \\ \hline
MemoryAgent & Extracción, grafo, indexado y resolución de conflictos. & Consolidar artefactos aceptados. \\ \hline
\end{tabular}
\caption{Herramientas usadas por agente.}
\end{center}
\end{table}

\section{Revisión humana y rutas de corrección}

La revisión humana forma parte del diseño. El Router llama a un puerto de
feedback. Ese puerto abstrae revisión por CLI, revisión por WebSocket y
autoaprobación en evals.

\begin{table}[H]
\begin{center}
\small
\begin{tabular}{|p{3cm}|p{4.8cm}|p{4.8cm}|} \hline
\textbf{Punto} & \textbf{Decisión humana} & \textbf{Motivo} \\ \hline
Investigación & Aprobar paradigmas o pedir uno adicional. & Evitar marcos científicos irrelevantes o incompletos. \\ \hline
Formalización & Seleccionar formulaciones concretas. & Reducir variantes a las que merece implementar. \\ \hline
Entorno & Proporcionar \texttt{env\_spec.json}. & Definir acciones, percepción y recompensas. \\ \hline
Razonamiento & Aprobar, rechazar o pedir nueva formalización. & Separar error matemático de error de adaptación. \\ \hline
Construcción & Aprobar modelos o pedir corrección. & Separar bug de código de mala especificación. \\ \hline
\end{tabular}
\caption{Puntos de revisión humana del pipeline.}
\end{center}
\end{table}

La revisión no solo permite continuar o parar. También redirige el pipeline al
nivel adecuado. Un error conceptual vuelve a Formalizer. Una mala adaptación
vuelve a Reasoner. Un bug de Python vuelve a Builder.

\begin{table}[H]
\begin{center}
\small
\begin{tabular}{|p{3.5cm}|p{4.4cm}|p{4.5cm}|} \hline
\textbf{Revisión} & \textbf{Acción posible} & \textbf{Efecto} \\ \hline
Investigación & Pedir un paradigma adicional. & Lanza un DeepResearcher nuevo. \\ \hline
Razonamiento & Rehacer formalización. & Repite Formalizer y después Reasoner. \\ \hline
Razonamiento & Rechazar una spec. & Reejecuta Reasoner para esa rama. \\ \hline
Construcción & Pedir nueva spec. & Regenera Reasoner y vuelve a Builder. \\ \hline
Construcción & Rechazar un modelo. & Reejecuta Builder solo para esa formulación. \\ \hline
\end{tabular}
\caption{Rutas de corrección activadas por revisión humana.}
\end{center}
\end{table}

\section{Artefactos, persistencia e integración}

El sistema prioriza la persistencia de artefactos. MinIO guarda contenido
completo. PostgreSQL guarda metadatos, ejecuciones, modelos aprobados y filas
de memoria. Esta separación evita guardar grandes bloques de texto o código en
SQL, pero mantiene referencias relacionales para listar runs y recuperar
estado.

\begin{verbatim}
research/{run_id}/
  report.md
  deep/{paradigm}.md
  formulations/{paradigm}.md
  env_spec.json
  pipeline_state.json
  trace.jsonl

models/{run_id}/
  reasoner/{paradigm}/{formulation}.json
  builder/{paradigm}/{formulation}_model.py
  builder/{paradigm}/test_{formulation}.py
  builder/{paradigm}/{formulation}_validation.json
\end{verbatim}

\begin{table}[H]
\begin{center}
\small
\begin{tabular}{|p{3.6cm}|p{4.3cm}|p{4.6cm}|} \hline
\textbf{Artefacto} & \textbf{Contenido} & \textbf{Uso} \\ \hline
\texttt{report.md} & Paradigmas, autores y conceptos. & Revisión de cobertura científica. \\ \hline
\texttt{deep/*.md} & Fundamentos, postulados, variables y referencias. & Base para formalización. \\ \hline
\texttt{formulations/*.md} & Variables, parámetros y ecuaciones. & Comparar alternativas implementables. \\ \hline
\texttt{env\_spec.json} & Acciones, percepciones y recompensas. & Fuente autoritativa del entorno. \\ \hline
\texttt{reasoner/*.json} & Reglas, mapeo y comportamientos esperados. & Entrada validable del Builder. \\ \hline
\texttt{builder/*.py} y tests & Código y pruebas. & Evidencia del contrato ejecutable. \\ \hline
\texttt{trace.jsonl} & Eventos, herramientas, tiempos y errores. & Auditoría y observabilidad. \\ \hline
\end{tabular}
\caption{Función de los artefactos generados.}
\end{center}
\end{table}

La reanudación comprueba que los artefactos referenciados por el estado existen
en MinIO. Si faltan formulaciones, specs o modelos, la reanudación se rechaza
para evitar continuar desde un estado incoherente.

\section{Sistema de memoria extendido}
\label{ap:memoria-extendida}

La memoria convierte el pipeline en un sistema acumulativo. Sin memoria, cada
ejecución empieza desde cero. Con memoria, los artefactos aceptados se
transforman en paradigmas, variables, postulados, formulaciones, parámetros,
modelos, relaciones y observaciones reutilizables.

La regla central es no guardar cualquier borrador generado por un LLM. Solo se
consolida memoria después de revisión humana o validación. Esta decisión reduce
ruido y evita que una formulación rechazada condicione ejecuciones futuras.

\subsection{Representaciones complementarias}

La memoria no usa un único almacén porque cada representación responde a una
pregunta distinta.

\begin{table}[H]
\begin{center}
\small
\begin{tabular}{|p{3.1cm}|p{4.2cm}|p{5.1cm}|} \hline
\textbf{Representación} & \textbf{Pregunta} & \textbf{Motivo} \\ \hline
Hechos dense y sparse & Qué significa un concepto o qué variables aparecen. & Dense recupera similitud semántica. Sparse recupera nombres, símbolos y autores exactos. \\ \hline
Grafo de conocimiento & De dónde sale una idea y cómo se relaciona con otras. & Permite procedencia mediante nodos y relaciones. \\ \hline
Metadatos relacionales & Qué está vigente, aceptado o supersedido. & Mantiene confianza, validez temporal y auditoría. \\ \hline
Observaciones & Qué se intentó antes y qué falló. & Evita repetir errores de adaptación o implementación. \\ \hline
\end{tabular}
\caption{Representaciones de la memoria.}
\end{center}
\end{table}

La búsqueda dense permite recuperar hechos aunque la consulta use términos
distintos. La búsqueda sparse con BM25 permite encontrar símbolos, slugs,
autores o identificadores exactos. El grafo permite reconstruir procedencia:
qué paper soporta un postulado, qué autor lo propuso, qué variables derivan de
él o qué modelo lo implementa.

\subsection{MemoryAgent}

El MemoryAgent no es un agente conversacional libre. Ejecuta un pipeline fijo
que usa modelos de lenguaje para extracción, pero mantiene control
determinista.

La extracción produce:

\begin{itemize}
\item nodos, como paradigmas, variables, papers, postulados, formulaciones y modelos;
\item relaciones, como soporta, implementa, usa variable o deriva de;
\item hechos y observaciones, como frases atómicas recuperables en ejecuciones futuras
\end{itemize}

Antes de escribir, el sistema normaliza identidades, poda nodos sin relación,
materializa relaciones estructurales y resuelve propuestas nuevas contra
paradigmas existentes. Este paso reduce duplicados semánticos.

\subsection{Riesgos de memoria}

\begin{table}[H]
\begin{center}
\small
\begin{tabular}{|p{3cm}|p{4.8cm}|p{4.8cm}|} \hline
\textbf{Riesgo} & \textbf{Efecto posible} & \textbf{Mitigación} \\ \hline
Duplicado semántico & El mismo paradigma aparece con nombres distintos. & Slugs canónicos, clasificación previa y revisión. \\ \hline
Hecho falso consolidado & Una alucinación afecta a runs futuros. & Escritura solo tras aceptación. \\ \hline
Recuperación irrelevante & El agente recibe contexto que distrae. & Búsqueda híbrida, reranking y filtros temporales. \\ \hline
Memoria obsoleta & Un hecho superado sigue apareciendo. & Vigencia temporal y supersesión en PostgreSQL. \\ \hline
Grafo incompleto & Se pierde procedencia de relaciones. & Artefactos completos en MinIO y snapshots de evals. \\ \hline
\end{tabular}
\caption{Riesgos de memoria y mitigaciones.}
\end{center}
\end{table}

\section{Desarrollo con agentes extendido}
\label{ap:desarrollo-agentes-extendido}

El desarrollo se organizó con agentes de programación y revisión humana. La
idea central fue separar decisiones de arquitectura y ejecución. El equipo
humano definía objetivo, restricciones y evidencia esperada. El agente
implementaba dentro de ese marco. Después se revisaban diffs, tests, logs,
artefactos y CI.

\subsection{Planificación con Linear}

Las tareas se planificaron en Linear. Cada tarea incluía objetivo, criterios de
aceptación, dependencias, prioridad, posible solapamiento y capacidad de
ejecutarse en paralelo.

Estos criterios fueron necesarios porque los agentes necesitan una definición
objetiva de terminado. Sin criterios explícitos, un agente puede generar código
funcional a primera vista pero incompleto desde producto o arquitectura.

\subsection{División entre Codex, Claude Code e Hydra}

El proyecto usó agentes distintos para tareas distintas:

\begin{itemize}
\item Codex se usó sobre todo para backend, infraestructura, persistencia, tests y razonamiento técnico;
\item Claude Code se usó sobre todo para frontend, interacción y experiencia de usuario;
\item Hydra se usó para ejecutar oleadas autónomas de tareas ya especificadas
\end{itemize}

Backend y frontend tienen ritmos de validación diferentes. En backend importan
contratos, persistencia, migraciones y límites entre fases. En frontend importan
coherencia visual, ergonomía, interacción y sensibilidad al detalle.

\subsection{Orquestación autónoma con Hydra}

Hydra no sustituía a Linear ni a la revisión humana. Convertía una lista de
tareas especificadas en oleadas de ejecución autónoma. Leía tareas pendientes,
agrupaba trabajo no conflictivo, lanzaba agentes en worktrees separados,
recogía evidencias y cerraba o creaba seguimiento.

\begin{table}[H]
\begin{center}
\small
\begin{tabular}{|p{3.2cm}|p{4.4cm}|p{4.8cm}|} \hline
\textbf{Paso} & \textbf{Responsable} & \textbf{Evidencia} \\ \hline
Descubrir tareas & Hydra leyendo Linear. & Lista de tareas, dependencias y prioridad. \\ \hline
Planificar oleada & Hydra. & Tareas ejecutables sin bloqueos ni solapamientos peligrosos. \\ \hline
Asignar modelo & Hydra y criterio humano previo. & Codex para backend; Claude Code para frontend. \\ \hline
Ejecutar tarea & Agente trabajador. & Worktree propio, implementación y pruebas locales. \\ \hline
Revisar salida & Agente trabajador. & Revisión de bugs, fallos silenciosos y convenciones. \\ \hline
Cerrar seguimiento & Hydra. & PR mergeada, nota o tarea nueva. \\ \hline
\end{tabular}
\caption{Fases del flujo autónomo con Hydra.}
\end{center}
\end{table}

Hydra mantenía su contexto pequeño. No cargaba todos los documentos, diffs y
logs de cada tarea. El agente trabajador leía el detalle dentro de su propia
sesión. Al terminar, devolvía resumen, ficheros tocados, estado de pruebas y
notas de seguimiento.

\subsection{Revisión de trabajo generado}

La revisión final de cada agente siguió 4 pasos:

\begin{itemize}
\item buscar bugs funcionales evidentes;
\item detectar fallos silenciosos, donde el sistema parece funcionar pero pierde datos;
\item revisar fallos de arquitectura, como acoplamientos indebidos;
\item comprobar convenciones del repositorio, estilo, tests y contratos locales
\end{itemize}

Este patrón fue necesario porque las tareas largas de desarrollo autónomo
siguen siendo difíciles. La metodología no consistió en confiar ciegamente en
un agente, sino en aislar tareas, exigir evidencia ejecutable y mantener un
humano responsable de producto y arquitectura.

\section{Validación extendida}
\label{ap:validacion-extendida}

\subsection{Integración continua}

El repositorio incluye un flujo de GitHub Actions llamado \texttt{CI}. Se
ejecuta en cada push a \texttt{main}, en cada pull request contra \texttt{main}
y de forma manual. También cancela ejecuciones antiguas de la misma rama cuando
llega un push nuevo.

La validación automática tiene 4 trabajos:

\begin{itemize}
\item \texttt{Lint + format}, que ejecuta \texttt{ruff} con \texttt{uv};
\item \texttt{Tests - shared}, que sincroniza dependencias y ejecuta tests comunes;
\item \texttt{Tests - phase1}, que ejecuta tests unitarios de fase 1;
\item \texttt{Tests - phase2}, que ejecuta tests unitarios de fase 2
\end{itemize}

Las pruebas excluyen marcas de integración, e2e y \texttt{real\_llm}, porque
esas suites dependen de servicios externos o ejecuciones caras. Las evals
online se lanzan explícitamente para cambios de más riesgo.

\subsection{Evals del sistema multiagente}

Las evals se definen como suites YAML. Cada suite declara etapas, entradas,
aprobaciones automáticas, presupuesto máximo y aserciones. Se ejecutan desde la
CLI y generan reportes Markdown y JSON.

\begin{verbatim}
uv run decisionlab eval run evals/suites/smoke.yaml
uv run decisionlab eval run evals/suites/full-pipeline.yaml
\end{verbatim}

El reporte incluye estado global, temas ejecutados, aserciones, crecimiento del
grafo, escrituras de memoria, duración por etapa, latencia de herramientas y
coste estimado.

\begin{table}[H]
\begin{center}
\small
\begin{tabular}{|p{3.2cm}|p{2.5cm}|p{6.6cm}|} \hline
\textbf{Suite} & \textbf{Etapas} & \textbf{Qué valida} \\ \hline
Smoke & Research. & Arranque básico y descubrimiento mínimo de paradigmas. \\ \hline
Full pipeline & Research a Build. & Investigación, formalización, razonamiento, generación, importación y llamada a \texttt{decide}. \\ \hline
\end{tabular}
\caption{Evals del flujo multiagente.}
\end{center}
\end{table}

Estas evals cubren fallos que no aparecen en pruebas unitarias aisladas. Por
ejemplo, comprueban que Researcher, Formalizer, Reasoner y Builder producen
artefactos compatibles entre sí.

\subsection{Evals de memoria}

Las evals de memoria comprueban Qdrant, Neo4j, PostgreSQL, canonicalización,
recuperación, deduplicación, crecimiento del grafo y validez temporal.

\begin{table}[H]
\begin{center}
\small
\begin{tabular}{|p{3.2cm}|p{2.5cm}|p{6.6cm}|} \hline
\textbf{Suite} & \textbf{Etapas} & \textbf{Qué valida} \\ \hline
Cumulative growth & Research. & Crecimiento controlado y reutilización de slugs. \\ \hline
Canonicalization & Research. & Reutilización de paradigmas conocidos. \\ \hline
Slug accuracy & Research. & Acierto contra oráculo de slugs y latencia de recuperación. \\ \hline
Memory retrieval & Research. & Recuperación de paradigmas presentes en memoria. \\ \hline
Retrieval quality & Offline. & Hechos sembrados aparecen en top-k. \\ \hline
Temporal correctness & Offline. & Validez por fecha y cadenas de supersesión. \\ \hline
\end{tabular}
\caption{Evals del sistema de memoria.}
\end{center}
\end{table}

Las aserciones combinan métricas agregadas y comprobaciones directas:

\begin{verbatim}
slug_hit_rate
kg_growth_rate
p95_below
memory_at_time
supersession_chain_length
retrieval_finds
\end{verbatim}

\subsection{Corpus cerrado de papers}

La evaluación más exigente se hizo con 2 casos basados en PDFs. En cada caso,
el sistema solo podía acceder a los papers incluidos en el ZIP de entrada. Las
herramientas de búsqueda devolvían metadatos y snippets de esos documentos.
Cuando el pipeline necesitaba leer un paper, la herramienta devolvía el texto
extraído completo.

Esta configuración mide si el pipeline puede descubrir paradigmas, proponer
formulaciones, razonar sobre un entorno, generar modelos y conservar artefactos
usando solo un corpus cerrado. Después se usó un juez LLM con la misma rúbrica
para ambos casos. El juez no participó en la generación.

\begin{table}[H]
\begin{center}
\small
\begin{tabular}{|p{1.6cm}|p{2.5cm}|p{2.2cm}|p{5.7cm}|} \hline
\textbf{Caso} & \textbf{Resultado} & \textbf{Tiempo / coste} & \textbf{Evidencia} \\ \hline
CASO1 & Pipeline PASS. Juez: 72/100 con reservas. & 6152.0 s / \$33.71 & 3 papers, 6 paradigmas, 18 specs, 15 modelos y KG de 433 nodos. \\ \hline
CASO2 & Pipeline PASS. Juez: 78/100 con reservas. & 4961.8 s / \$24.47 & 3 papers, 4 paradigmas, 12 specs, 12 modelos y KG de 272 nodos. \\ \hline
\end{tabular}
\caption{Resultados de evals con corpus cerrado.}
\end{center}
\end{table}

El veredicto aprobado con reservas debe leerse como señal de utilidad, no como
certificación científica final. La señal positiva es que el sistema produjo
artefactos revisables con corpus cerrado. La señal negativa es que el juez
encontró problemas reales de taxonomía, registro y trazabilidad.

\begin{table}[H]
\begin{center}
\scriptsize
\begin{tabular}{|p{2.3cm}|p{1.5cm}|p{3cm}|p{3cm}|p{2.2cm}|} \hline
\textbf{Reserva} & \textbf{Severidad} & \textbf{Efecto} & \textbf{Mitigación} & \textbf{Estado} \\ \hline
Spec inválida en CASO1 & Alta & No debe generar un modelo aceptado. & Reasoner marca salida inválida y conserva problemas. & Revisión experta. \\ \hline
Registro DDM desajustado & Alta & Puede apuntar a clase o test distinto. & Bundle conserva rutas y filas para auditar. & Detectado por juez. \\ \hline
Errores en KG & Media & Reducen trazabilidad automática. & Artefactos, snapshots y conteos permiten localizar fallos. & Limitación conocida. \\ \hline
Solapamiento taxonómico & Media & Puede mezclar paradigmas relacionados. & Revisión humana de paradigmas y formulaciones. & Necesita experto. \\ \hline
Duplicados de ficheros & Baja & Dificultan revisión manual. & Tabla de modelos conserva ruta canónica. & Higiene pendiente. \\ \hline
\end{tabular}
\caption{Reservas detectadas por el juez.}
\end{center}
\end{table}

La evidencia versionada queda en los bundles y reportes siguientes:

\begin{small}
\begin{itemize}
\item bundle CASO1, \path{evals/reports/2026-06-29-caso1-pdf-corpus};
\item bundle CASO2, \path{evals/reports/2026-06-29-caso2-pdf-corpus};
\item reportes del juez en \path{evals/judge-reports};
\item snapshot de infraestructura en \path{evals/infra-snapshots}
\end{itemize}
\end{small}

\section{Limitaciones y trabajo futuro extendido}
\label{ap:limitaciones-extendidas}

La principal limitación es que la validez científica final de un modelo no
queda garantizada por compilar o pasar tests. Las pruebas automáticas verifican
propiedades funcionales, pero la adecuación psicológica o neurocientífica
requiere revisión experta.

Otra limitación es la dependencia de servicios externos y modelos de lenguaje.
Las ejecuciones completas tienen coste, latencia y variabilidad. Por eso las
evals se interpretan como señales de regresión interna, no como benchmark
externo definitivo.

También hay margen de mejora en la memoria. Aunque el sistema distingue hechos,
relaciones y observaciones, la calidad de extracción depende de los artefactos
producidos y de los criterios de consolidación. La canonicalización de
paradigmas, detección de contradicciones y calibración de confianza pueden
mejorar con más ejecuciones.

Las líneas futuras principales son:

\begin{itemize}
\item ampliar el conjunto de paradigmas y ejecutar más casos con corpus cerrado;
\item validar resultados con expertos del dominio;
\item comparar modelos generados con implementaciones manuales;
\item estudiar si los modelos reproducen patrones esperados en simulaciones;
\item mejorar la extracción de observaciones procedurales;
\item usar más el grafo para consultas multi-hop sobre procedencia científica;
\item reforzar la integración con fase 2 y retroalimentar memoria desde resultados de simulación
\end{itemize}
